\chapter{Einleitung}
Licht besitzt laut Welle-Teilchen-Dualismus einen Teilchen- und einen Wellencharakter \cite{meschede2006gerthsen}.
Im Wellenbild ist die Wellenlänge, und damit auch die Frequenz, ausschlaggebend dafür ob das Licht sichtbar ist und welche Farbe es hat.
Zum beispiel ist das Licht zwischen einer Wellenlänge von \SIrange{640}{430}{\nano\meter} im sichtbaren bereich. 
Aufgrund der elektromagnetischen Eigenschaften des Lichts wird das gesamte Spektrum auch als elektromagnetisches Spektrum, 
kurz EM-Spektrum, bezeichnet. 
Dieses ist in \autoref{fig:Spek} dargestellt. \\ 
\begin{figure}[H]
    \centering
    \includegraphics[width=0.8\textwidth]{content/bilder/Em-Spek.png}
    \caption{Elektromagnetisches-Spektrum mit Frequenzskala und typischen anwendungsbereichen. 
    Rot Markiert ist der THz-Bereich, welcher zwischen dem Mikrowellen und dem Sichtbaren Spektrum liegt. 
    Charakteristische Kenngrößen der THz-Strahlung sind darunter angegeben \cite{Williams_2006}.}
    \label{fig:Spek}
\end{figure}
Wobei ein Großteil dieses Spektrum ohne viel aufwand zu erzeugen ist, gibt es bei gewissen Frequenzbereichen erhebliche Herausforderungen.
Dieser herausfordernde Bereich ist die sogenannte Teraheartz-Lücke, sie beschreibt 
einen Frequenzbereich in welchem es technisch besonders schwierig ist elektromagnetische-Wellen zu erzeugen und zu detektieren. 
Das ist in \autoref{fig:THz-lucke} zu sehen. 
Dort ist zu sehen, dass auf der seiter der niedrigeren Frequenzen als THz also auf der Linken seite die Leistung von Elektronischen erzeugern 
bei verschiedenen Frequenzen aufgetragen ist. Auf der Rechten seite sind die Leistungen von Quantenkaskaden-Lasern bei verschiedenen Frequenzen 
auufgetragen. Auf der Abbildung ist somit zu sehen, dass knapp über \SI{1000}{GHz} also \SI{1}{THz} die Leistungen von beiden seiten 
nahe 0 gehen. \\
Diese THz-Lücke liegt im Frequenzbereich von \SIrange{0.1}{10}{\tera\hertz} oder, im Wellenlängen spektrum von \SIrange{3000}{30}{\micro\meter} \cite{Williams_2006}.
Frequenzen unterhalb der THz-Lücke sind gut durch elektronische Oszillatoren zu erzeugen. 
Bei ansteigenden Frequenzen entstehen aber immer mehr Herausforderungen, zum Beispiel durch 
parasitäre Kapazitäten, welche zur Folge haben, dass immer weniger Leistung als Strahlung 
abgegeben werden kann.  
Dazu können einfache Schwarzkörper oder Diodenlaser hohe Leistung 
im Bereich über der THz-Lücke erzeugen. 
Dies nimmt aber stark mit sinkender Frequenz ab.
Somit ergibt sich eine Lücke im Spektrum, in welcher es sehr schwierig ist, eine elektromagnetische Welle mit THz-Frequenzen
und hoher Leistung zu erzeugen. 
\begin{figure}[H]
    \centering
    \includegraphics[width=0.8\textwidth]{content/bilder/THz-gap.png}
    \caption{Graf der THz-Lücke Wobei Links der Lücke bei niedrigeren Frequenzen die Leistung von Elektrischen Strahlungserzeugern zu sehen ist. 
    Auf der Rechten seite sind die Leistung von Quantenkaskaden-Laser eingezeitnet.
    Die beiden Leistungen sinken in richtung der eingzeichneten THz-Lücke in der mitte \cite{thz_gap}.}
    \label{fig:THz-lucke}
\end{figure}
Dabei ist THz-Strahlung aus verschiedenen Gründen interessant. THz-Strahlung kann wie 
Röntgenstrahlung zum Finden und zur Detektion von Tumoren und anderen Gewebeanomalien genutzt werden, wobei es 
nicht die Erbschäden wie Röntgenstrahlung verursacht \cite{pickwell2006biomedical}. 
In den Naturwissenschaften wird sie zur nicht-destruktiven Materialbestimmung und Untersuchung, 
sowie nieder energetischen Anregung benutzt. \cite{PAWAR2013157}.
Die Erzeugung von Strahlung innerhalb der THz-Lücke ist heute durch die Nutzung von verschiedenen THz-Generations-Techniken, wie zum Beispiel photoleitenden Antennen 
\cite{10.1117/1.OE.56.1.010901}, optischer Gleichrichtung und THz-Quantenkaskaden \cite{Dean_2014}, 
in höheren Intensitäten möglich.
In dieser Bachelorarbeit liegt besonderes Interesse auf der optischen Gleichrichtung, da 
diese der Hauptmechanismus, mit welchem die organischen Emitter THz-Strahlung erzeugen, ist. \\ \\
Sie werden mit anorganischen Kristallen verglichen. ZnTe ist ein anorganischer Kristall der ebenfalls durch optische Gleichrichtung THz-Strahlung erzeugen kann.
Durch zu hohe optische Leistung auf dem Kristall kann dieser zerstört werden, 
durch die geringe Anzahl an Produzenten und der ökonomischen Entwicklung der letzten Jahre hat sich der Preis hier in den letzten 
fünf Jahren verdoppelt. \\ \\
In dieser Bachelorarbeit werden zwei organische Kristallsysteme auf ihre Eigenschaften bezüglich der THz-Generation untersucht und mit 
Generierung in ZnTe verglichen. \\
Die organischen Kristalle werden in der Arbeitsgruppe $Nachgucken welche gruppe das war$an der TU Dortmund gezüchtet/hergestellt, 
welches eine schnelle Iteration bezüglich des Designs der Kristalle ermöglicht.  \\
So können Anwendungsspezifische Kristalle hinsichtlich Größe und Stärke gezüchtet werden, 
im Kontrast zu kommerziell erwerbbaren anorganischen Kristallen. \\
Anschließend wird diskutiert, ob diese Kristalle die ZnTe-Kristalle ersetzen können. 
Die vorliegende Arbeit beginnt in Kapitel zwei mit einer theoretischen Beschreibung der Generation und Detektion von 
THz-Strahlung um darauf in Kapitel drei aufbauend die Terahertz-Zeitbereichsspektroskopie zu erklären. 
Diese wird folgend  mit THz-TDS abgekürzt, aus der englischen Bezeichnung Terahertz-time domain spectroscopy.
Die mit dieser Messmethode gemessenen Ergebnisse der verschiedenen Kristalle werden in Kapitel vier dargestellt und diskutiert.
Kapitel fünf schließt mit einem Überblick über die Ergebnisse und einem Ausblick in weitere Messungen ab.